\section{Analyse de \texttt{full\_soft/}}

\subsection{Organisation generale}

Lors de l'audit initial, le miroir historique \path{full_soft/} etait organise en plusieurs sous-ensembles. Ces chemins sont conserves ici comme references documentaires, mais ils ne correspondent plus a des fichiers presents dans \path{main}.

\begin{longtable}{>{\raggedright\arraybackslash}p{0.28\linewidth} >{\raggedright\arraybackslash}p{0.62\linewidth}}
\textbf{Emplacement} & \textbf{Role observe} \\
\hline
\path{full_soft/code} & Code Python principal de controle vehicule. \\
\path{full_soft/code/algorithm} & Algorithmes de navigation, filtrage LiDAR, controle direction/vitesse et strategies. \\
\path{full_soft/code/Interfaces} & Interfaces materielles et simulees pour LiDAR, moteur, direction, camera, serie et abstractions. \\
\path{full_soft/Simulateur} & Simulateur Python historique, generation de pistes et logs. \\
\path{full_soft/scripts} & Scripts d'installation, configuration PWM et udev. \\
\path{full_soft/docs} & Documentation Raspberry Pi, WiFi, VNC, VSCode, hardware et logiciel. \\
\path{full_soft/documentation} & Inventaire et documents d'exigences/tests. \\
\end{longtable}

\subsection{Points d'entree}

Deux points d'entree principaux coexistent:

\begin{itemize}
  \item \path{full_soft/code/main.py}, qui charge \path{new-config.json}, instancie les interfaces, cree \path{VoitureAlgorithm} et lance une boucle continue;
  \item \path{full_soft/code/main_refactored.py}, qui charge \path{config.json}, initialise le logging, les interfaces, \path{VehicleAlgorithm_v1} et propose aussi un mode de simulation.
\end{itemize}

Cette coexistence etait un fait confirme lors de l'audit du miroir historique. Elle indique une transition entre une version initiale et une version refactorisee, mais aucun fichier ne designait explicitement l'un des deux comme unique point d'entree officiel.

\subsection{Interfaces}

Le dossier \path{Interfaces} contient des abstractions et implementations pour:

\begin{itemize}
  \item le LiDAR reel ou simule;
  \item le moteur;
  \item la direction;
  \item la camera;
  \item la liaison serie;
  \item des versions factices ou simulees.
\end{itemize}

Cette separation est positive: elle permet de remplacer une implementation materielle par une implementation de simulation ou de test. Elle reste cependant specifique a cette pile Python et n'expose pas nativement de topics ROS.
